\documentclass[aspectratio=169]{beamer}

\usepackage[utf8]{inputenc}
\usepackage[T1]{fontenc}
\usepackage[spanish]{babel}
\usepackage{amsmath}
\usepackage{amssymb}
\usepackage{booktabs}
\usepackage{graphicx}
\usepackage{tabularx}
\usepackage{ragged2e}

\usetheme{Madrid}
\usecolortheme{default}

\definecolor{untrmblue}{RGB}{18,60,105}
\definecolor{untrmgreen}{RGB}{28,122,92}
\definecolor{softgray}{RGB}{245,247,250}
\definecolor{accentgold}{RGB}{214,157,49}

\setbeamercolor{structure}{fg=untrmblue}
\setbeamercolor{palette primary}{bg=untrmblue,fg=white}
\setbeamercolor{palette secondary}{bg=untrmgreen,fg=white}
\setbeamercolor{palette tertiary}{bg=accentgold,fg=white}
\setbeamercolor{frametitle}{bg=untrmblue,fg=white}
\setbeamercolor{title}{fg=white,bg=untrmblue}
\setbeamercolor{block title}{bg=untrmblue,fg=white}
\setbeamercolor{block body}{bg=softgray,fg=black}
\setbeamertemplate{navigation symbols}{}
\setbeamertemplate{footline}[frame number]
\setbeamerfont{author}{size=\small}
\setbeamerfont{institute}{size=\scriptsize}
\setbeamersize{text margin left=8mm,text margin right=8mm}
\setbeamertemplate{itemize item}{\color{untrmgreen}$\blacktriangleright$}
\setbeamertemplate{itemize subitem}{\color{accentgold}$\bullet$}

\graphicspath{{outputs/}}

\title[Modelo Mundial 2026]{Modelo Mundial 2026}
\subtitle{Predicción de resultados con modelos de regresión y contexto competitivo}
\author[Equipo de trabajo]{
    Bustamante Vela Jean Frank\\
    Marreros Cortegana Miguel Angel\\
    Fernandez Campos Valentin\\
    Salon Trigoso Frank\\
    Tuesta Chuquizuta Maria Carmen
}
\institute[UNTRM]{
    Universidad Nacional Toribio Rodriguez de Mendoza de Amazonas\\
    Curso: Modelos de Regresión y Aprendizaje Supervisado
}
\date{\today}

\begin{document}

%------------------------------------------------
\begin{frame}
    \titlepage
\end{frame}

%------------------------------------------------
\begin{frame}{Idea central del proyecto}
\justifying
El proyecto busca estimar los posibles resultados de partidos del Mundial 2026 a partir de información histórica y variables deportivas medibles.

\vspace{0.25cm}
\begin{block}{Pregunta guía}
¿Qué tan probable es que un equipo gane, empate o pierda un partido, y qué marcador aparece como más probable según los modelos?
\end{block}

\vspace{0.2cm}
La propuesta combina modelos de regresión, probabilidades de marcadores y un ajuste por contexto competitivo para simular escenarios de fase de grupos.
\end{frame}

%------------------------------------------------
\begin{frame}{Objetivo del estudio}
\begin{block}{Objetivo general}
Construir un sistema predictivo que estime goles esperados, probabilidades de resultado y clasificación a la siguiente ronda del Mundial 2026.
\end{block}

\vspace{0.2cm}
\begin{itemize}
    \item Estimar goles esperados del equipo local y visitante.
    \item Comparar diferentes modelos de regresión supervisada.
    \item Transformar goles esperados en probabilidades de victoria, empate y derrota.
    \item Simular tablas de grupos, mejores terceros y cruces de 16avos.
\end{itemize}
\end{frame}

%------------------------------------------------
\begin{frame}{Información recolectada}
\justifying
Para representar la fuerza deportiva de cada selección se integraron distintas fuentes de información.

\vspace{0.25cm}
\begin{columns}[T]
\column{0.5\textwidth}
\begin{block}{Datos deportivos}
\begin{itemize}
    \item Resultados históricos internacionales.
    \item Goles anotados y recibidos recientemente.
    \item Puntos obtenidos en partidos previos.
    \item Historial de enfrentamientos directos.
\end{itemize}
\end{block}

\column{0.5\textwidth}
\begin{block}{Datos de contexto}
\begin{itemize}
    \item Ranking FIFA.
    \item Rating Elo.
    \item Condición de campo neutral.
    \item Ventaja de localía.
    \item Información de calendario y grupos.
\end{itemize}
\end{block}
\end{columns}

\vspace{0.2cm}
La base final de entrenamiento contiene \textbf{5 901 partidos} y el fixture analizado contiene \textbf{72 encuentros}.
\end{frame}

%------------------------------------------------
\begin{frame}{Variables predictoras}
\justifying
El modelo utiliza 13 variables principales, agrupadas en cuatro bloques.

\vspace{0.2cm}
\begin{tabularx}{\textwidth}{>{\bfseries}p{0.23\textwidth}X}
\toprule
Bloque & Variables consideradas \\
\midrule
Forma reciente local & Goles a favor, goles en contra, puntos y partidos recientes del equipo local.\\
Forma reciente visitante & Goles a favor, goles en contra, puntos y partidos recientes del equipo visitante.\\
Fuerza comparativa & Diferencia en ranking FIFA, diferencia en rating Elo e historial directo.\\
Contexto del partido & Campo neutral y ventaja de localía.\\
\bottomrule
\end{tabularx}

\vspace{0.3cm}
Estas variables permiten representar rendimiento, fortaleza relativa y condiciones del encuentro.
\end{frame}

%------------------------------------------------
\begin{frame}{Vector matemático de entrada}
\justifying
Cada partido se representa mediante un vector de 13 características:

\[
X = (X_1, X_2, X_3, \dots, X_{13})
\]

\vspace{0.15cm}
\begin{block}{Interpretación}
El vector $X$ resume la información deportiva y contextual disponible antes del partido: forma reciente, goles, puntos, ranking, rating Elo, historial directo y localía.
\end{block}

\vspace{0.2cm}
Sobre este vector se estiman dos respuestas numéricas: los goles esperados del equipo local y los goles esperados del equipo visitante.
\end{frame}

%------------------------------------------------
\begin{frame}{Correspondencia de variables}
\scriptsize
\centering
\resizebox{\textwidth}{!}{%
\begin{tabular}{clcl}
\toprule
Variable & Significado & Variable & Significado \\
\midrule
$X_1$ & Goles a favor local & $X_8$ & Partidos recientes visitante \\
$X_2$ & Goles en contra local & $X_9$ & Diferencia ranking FIFA \\
$X_3$ & Puntos recientes local & $X_{10}$ & Diferencia rating Elo \\
$X_4$ & Partidos recientes local & $X_{11}$ & Historial directo \\
$X_5$ & Goles a favor visitante & $X_{12}$ & Campo neutral \\
$X_6$ & Goles en contra visitante & $X_{13}$ & Ventaja de localía \\
$X_7$ & Puntos recientes visitante & & \\
\bottomrule
\end{tabular}
}

\vspace{0.35cm}
\begin{block}{Nota}
Los coeficientes numéricos presentados se interpretan sobre variables estandarizadas.
\end{block}
\end{frame}

%------------------------------------------------
\begin{frame}{Variables objetivo}
\justifying
El problema se formuló como una predicción numérica de goles.

\vspace{0.3cm}
\begin{columns}[T]
\column{0.5\textwidth}
\begin{block}{Goles del local}
\[
y_{\text{home}} = \text{goles anotados por el equipo local}
\]
\end{block}

\column{0.5\textwidth}
\begin{block}{Goles del visitante}
\[
y_{\text{away}} = \text{goles anotados por el equipo visitante}
\]
\end{block}
\end{columns}

\vspace{0.3cm}
Con estas dos estimaciones se obtiene una expectativa ofensiva para cada selección y luego se construyen probabilidades de marcador.
\end{frame}

%------------------------------------------------
\begin{frame}{Modelos evaluados}
\small
\begin{tabularx}{\textwidth}{>{\bfseries}p{0.28\textwidth}X}
\toprule
Modelo & Aporte dentro del análisis \\
\midrule
Regresión Lineal & Estima goles mediante una relación lineal entre variables y resultado.\\
Ridge & Regresión lineal regularizada; reduce sobreajuste y estabiliza coeficientes.\\
Random Forest & Ensamble de árboles que captura relaciones no lineales e interacciones.\\
Gradient Boosting & Ensamble secuencial que corrige errores de modelos anteriores.\\
Poisson Regressor & Modelo adecuado para conteos, porque los goles son valores discretos no negativos.\\
\bottomrule
\end{tabularx}

\vspace{0.25cm}
La comparación permitió seleccionar el modelo con mejor equilibrio entre error e interpretación.
\end{frame}

%------------------------------------------------
\begin{frame}{Regresión Lineal Múltiple}
\justifying
La Regresión Lineal modela la cantidad esperada de goles como una combinación ponderada de las variables del partido.

\[
\hat{y} = \beta_0 + \beta_1X_1 + \beta_2X_2 + \cdots + \beta_{13}X_{13}
\]

\vspace{0.2cm}
Como se predicen goles de ambos equipos, se plantean dos ecuaciones:

\[
\hat{y}_{\text{home}} = \beta_{0,h} + \sum_{i=1}^{13}\beta_{i,h}X_i
\]

\[
\hat{y}_{\text{away}} = \beta_{0,a} + \sum_{i=1}^{13}\beta_{i,a}X_i
\]
\end{frame}

%------------------------------------------------
\begin{frame}[t]{Coeficientes estimados: Regresión Lineal}
\scriptsize
\justifying
Coeficientes sobre variables estandarizadas. El signo indica si la variable aumenta o reduce la expectativa de goles.

\vspace{0.15cm}
\begin{columns}[T]
\column{0.5\textwidth}
\begin{block}{Predicción local}
\centering $\beta_0 = 1.6249$
\vspace{0.05cm}
\resizebox{\linewidth}{!}{%
\begin{tabular}{lr}
\toprule
Variable & $\beta$ \\
\midrule
Diferencia FIFA & 0.6405 \\
Goles recibidos visitante & 0.4051 \\
Goles a favor local & 0.2461 \\
Puntos visitante & 0.1892 \\
Puntos local & -0.1613 \\
\bottomrule
\end{tabular}
}
\end{block}

\column{0.5\textwidth}
\begin{block}{Predicción visitante}
\centering $\beta_0 = 1.1890$
\vspace{0.05cm}
\resizebox{\linewidth}{!}{%
\begin{tabular}{lr}
\toprule
Variable & $\beta$ \\
\midrule
Diferencia FIFA & -0.4475 \\
Goles recibidos local & 0.3959 \\
Goles a favor visitante & 0.2111 \\
Puntos local & 0.1464 \\
Puntos visitante & -0.1452 \\
\bottomrule
\end{tabular}
}
\end{block}
\end{columns}
\end{frame}

%------------------------------------------------
\begin{frame}[t]{Regresión Lineal con valores del proyecto}
\scriptsize
\justifying
Después de estimar los coeficientes, el modelo queda expresado de forma aproximada con los pesos principales:

\vspace{0.1cm}
\begin{block}{Goles esperados del local}
\[
\begin{aligned}
\hat{y}_{\text{home}}
&\approx
1.6249+0.6405X_9+0.4051X_6\\
&\quad +0.2461X_1+0.1892X_7-0.1613X_3+\cdots
\end{aligned}
\]
\end{block}

\begin{block}{Goles esperados del visitante}
\[
\begin{aligned}
\hat{y}_{\text{away}}
&\approx
1.1890-0.4475X_9+0.3959X_2\\
&\quad +0.2111X_5+0.1464X_3-0.1452X_7+\cdots
\end{aligned}
\]
\end{block}

\vspace{0.1cm}
En esta forma se observa qué variables tienen mayor influencia en la predicción de goles.
\end{frame}

%------------------------------------------------
\begin{frame}{Regresión Ridge}
\justifying
Ridge mantiene la estructura lineal, pero agrega regularización $L_2$ para controlar el tamaño de los coeficientes y reducir el sobreajuste.

\[
\min_{\beta}
\left[
\sum_{i=1}^{n}(y_i-\hat{y}_i)^2
+
\alpha\sum_{j=1}^{13}\beta_j^2
\right]
\]

\vspace{0.25cm}
\begin{block}{Razón de uso}
Es útil cuando existen variables relacionadas entre sí, como ranking FIFA, rating Elo, goles recientes y puntos recientes. Por eso fue elegido como modelo final.
\end{block}
\end{frame}

%------------------------------------------------
\begin{frame}[t]{Coeficientes del modelo final Ridge}
\scriptsize
\justifying
Coeficientes principales del modelo seleccionado. La diferencia FIFA y la fortaleza reciente tienen mayor peso.

\vspace{0.15cm}
\begin{columns}[T]
\column{0.5\textwidth}
\begin{block}{Goles esperados del local}
\centering $\beta_0 = 1.6249$
\vspace{0.05cm}
\resizebox{\linewidth}{!}{%
\begin{tabular}{lr}
\toprule
Variable & $\beta$ \\
\midrule
Diferencia FIFA & 0.6397 \\
Goles recibidos visitante & 0.4044 \\
Goles a favor local & 0.2456 \\
Puntos visitante & 0.1880 \\
Puntos local & -0.1603 \\
\bottomrule
\end{tabular}
}
\end{block}

\column{0.5\textwidth}
\begin{block}{Goles esperados del visitante}
\centering $\beta_0 = 1.1890$
\vspace{0.05cm}
\resizebox{\linewidth}{!}{%
\begin{tabular}{lr}
\toprule
Variable & $\beta$ \\
\midrule
Diferencia FIFA & -0.4469 \\
Goles recibidos local & 0.3953 \\
Goles a favor visitante & 0.2106 \\
Puntos local & 0.1454 \\
Puntos visitante & -0.1443 \\
\bottomrule
\end{tabular}
}
\end{block}
\end{columns}
\end{frame}

%------------------------------------------------
\begin{frame}[t]{Ridge con valores del proyecto}
\scriptsize
\justifying
Ridge fue el modelo seleccionado. Su forma reemplazada mantiene coeficientes similares a la regresión lineal, pero regularizados.

\vspace{0.1cm}
\begin{block}{Modelo final para goles del local}
\[
\begin{aligned}
\hat{y}_{\text{home}}
&\approx
1.6249+0.6397X_9+0.4044X_6\\
&\quad +0.2456X_1+0.1880X_7-0.1603X_3+\cdots
\end{aligned}
\]
\end{block}

\begin{block}{Modelo final para goles del visitante}
\[
\begin{aligned}
\hat{y}_{\text{away}}
&\approx
1.1890-0.4469X_9+0.3953X_2\\
&\quad +0.2106X_5+0.1454X_3-0.1443X_7+\cdots
\end{aligned}
\]
\end{block}
\end{frame}

%------------------------------------------------
\begin{frame}{Modelos de ensamble}
\small
\begin{block}{Random Forest}
Promedia la predicción de múltiples árboles de decisión:
\[
\hat{y} = \frac{1}{B}\sum_{b=1}^{B}T_b(X)
\]
\end{block}

\vspace{0.15cm}
\begin{block}{Gradient Boosting}
Construye modelos secuenciales donde cada nuevo árbol corrige parte del error anterior:
\[
F_m(X)=F_{m-1}(X)+\nu h_m(X)
\]
\[
r_{im}=y_i-F_{m-1}(X_i)
\]
\end{block}
\end{frame}

%------------------------------------------------
\begin{frame}[t]{Importancia de variables en modelos de ensamble}
\scriptsize
\justifying
Random Forest y Gradient Boosting no generan coeficientes $\beta$ como los modelos lineales. En su lugar, se analizan importancias de variables.

\vspace{0.15cm}
\begin{columns}[T]
\column{0.5\textwidth}
\begin{block}{Random Forest}
\resizebox{\linewidth}{!}{%
\begin{tabular}{lrr}
\toprule
Variable & Local & Visitante \\
\midrule
Diferencia FIFA & 0.5439 & 0.4807 \\
Goles recibidos rival & 0.1283 & 0.1831 \\
Goles a favor propio & 0.0785 & 0.0789 \\
Puntos recientes & 0.0449 & 0.0472 \\
Historial directo & -- & 0.0454 \\
\bottomrule
\end{tabular}
}
\end{block}

\column{0.5\textwidth}
\begin{block}{Gradient Boosting}
\resizebox{\linewidth}{!}{%
\begin{tabular}{lrr}
\toprule
Variable & Local & Visitante \\
\midrule
Diferencia FIFA & 0.5033 & 0.4399 \\
Goles recibidos rival & 0.1571 & 0.1952 \\
Goles a favor propio & 0.0771 & 0.0741 \\
Puntos recientes & 0.0471 & 0.0579 \\
Historial directo & 0.0412 & 0.0566 \\
\bottomrule
\end{tabular}
}
\end{block}
\end{columns}
\end{frame}

%------------------------------------------------
\begin{frame}[t]{Ensamble con valores obtenidos}
\scriptsize
\justifying
En estos modelos no se reemplaza una ecuación lineal con betas, porque la predicción se construye con árboles. Por ello se reportan salidas e importancias.

\vspace{0.15cm}
\begin{columns}[T]
\column{0.5\textwidth}
\begin{block}{Random Forest}
\[
\hat{y}_{\text{home}}=1.345,\qquad
\hat{y}_{\text{away}}=1.265
\]
\[
MAE_{\text{prom}}=0.9852,\quad RMSE_{\text{prom}}=1.3162
\]
Variable más influyente: diferencia FIFA.
\end{block}

\column{0.5\textwidth}
\begin{block}{Gradient Boosting}
\[
\hat{y}_{\text{home}}=1.259,\qquad
\hat{y}_{\text{away}}=1.143
\]
\[
MAE_{\text{prom}}=0.9939,\quad RMSE_{\text{prom}}=1.3434
\]
Mayor exactitud 1X2: $0.6274$.
\end{block}
\end{columns}

\vspace{0.15cm}
Los valores de $\hat{y}_{\text{home}}$ y $\hat{y}_{\text{away}}$ corresponden al partido Egypt vs Iran.
\end{frame}

%------------------------------------------------
\begin{frame}{Modelo de Poisson para goles}
\justifying
Los goles son datos de conteo: $0,1,2,3,\dots$. Por ello se considera una distribución de Poisson con parámetro $\lambda$.

\[
Y \sim \operatorname{Poisson}(\lambda)
\]

\[
\log(\lambda)=\beta_0+\beta_1X_1+\cdots+\beta_{13}X_{13}
\]

\[
\lambda=\exp\left(\beta_0+\sum_{i=1}^{13}\beta_iX_i\right)
\]

\vspace{0.1cm}
Así se obtienen $\lambda_{\text{home}}$ y $\lambda_{\text{away}}$, que representan los goles esperados de cada equipo.
\end{frame}

%------------------------------------------------
\begin{frame}[t]{Parámetros estimados del modelo Poisson}
\scriptsize
\justifying
En Poisson los coeficientes actúan sobre $\log(\lambda)$; un coeficiente positivo incrementa multiplicativamente la tasa esperada.

\vspace{0.15cm}
\begin{columns}[T]
\column{0.5\textwidth}
\begin{block}{Tasa del local}
\centering $\beta_0 = 0.3431$
\vspace{0.05cm}
\resizebox{\linewidth}{!}{%
\begin{tabular}{lr}
\toprule
Variable & $\beta$ \\
\midrule
Diferencia FIFA & 0.3990 \\
Goles recibidos visitante & 0.1880 \\
Goles a favor local & 0.1341 \\
Puntos local & -0.0923 \\
Historial directo & 0.0786 \\
\bottomrule
\end{tabular}
}
\end{block}

\column{0.5\textwidth}
\begin{block}{Tasa del visitante}
\centering $\beta_0 = 0.0073$
\vspace{0.05cm}
\resizebox{\linewidth}{!}{%
\begin{tabular}{lr}
\toprule
Variable & $\beta$ \\
\midrule
Diferencia FIFA & -0.3717 \\
Goles recibidos local & 0.1966 \\
Goles a favor visitante & 0.1448 \\
Historial directo & -0.0996 \\
Puntos visitante & -0.0933 \\
\bottomrule
\end{tabular}
}
\end{block}
\end{columns}
\end{frame}

%------------------------------------------------
\begin{frame}[t]{Poisson con valores del proyecto}
\scriptsize
\justifying
En Poisson se reemplaza la ecuación de $\log(\lambda)$ con los coeficientes estimados.

\vspace{0.1cm}
\begin{block}{Tasa esperada del local}
\[
\begin{aligned}
\log(\lambda_{\text{home}})
&\approx
0.3431+0.3990X_9+0.1880X_6\\
&\quad +0.1341X_1-0.0923X_3+0.0786X_{11}+\cdots
\end{aligned}
\]
\[
\lambda_{\text{home}}=\exp(\log(\lambda_{\text{home}}))
\]
\end{block}

\begin{block}{Tasa esperada del visitante}
\[
\begin{aligned}
\log(\lambda_{\text{away}})
&\approx
0.0073-0.3717X_9+0.1966X_2\\
&\quad +0.1448X_5-0.0996X_{11}-0.0933X_7+\cdots
\end{aligned}
\]
\end{block}
\end{frame}

%------------------------------------------------
\begin{frame}{Conversión de goles a probabilidades}
\justifying
Los goles esperados se convierten en tasas de anotación:

\[
\lambda_{\text{home}}, \quad \lambda_{\text{away}}
\]

Luego se calcula la probabilidad de cada marcador posible:

\[
P(\text{Home}=a)=
\frac{e^{-\lambda_{\text{home}}}\lambda_{\text{home}}^a}{a!}
\]

\[
P(\text{Away}=b)=
\frac{e^{-\lambda_{\text{away}}}\lambda_{\text{away}}^b}{b!}
\]

\[
P(a,b)=P(\text{Home}=a)\times P(\text{Away}=b)
\]

\vspace{0.2cm}
A partir de la matriz de marcadores se obtienen tres escenarios:

\begin{itemize}
    \item Victoria local: marcadores donde el local anota más.
    \item Empate: marcadores con igual número de goles.
    \item Victoria visitante: marcadores donde el visitante anota más.
\end{itemize}
\end{frame}

%------------------------------------------------
\begin{frame}[t]{Ejemplo numérico de $\lambda$: Egypt vs Iran}
\scriptsize
\centering
Para el partido analizado, cada modelo produjo tasas de goles esperados y probabilidades 1X2.

\vspace{0.2cm}
\resizebox{\textwidth}{!}{%
\begin{tabular}{lccccc}
\toprule
Modelo & $\lambda_{\text{home}}$ & $\lambda_{\text{away}}$ & Local & Empate & Visitante \\
\midrule
Ridge & \textbf{1.119} & \textbf{1.521} & 26.960\% & 27.481\% & \textbf{45.559\%} \\
Regresión Lineal & 1.118 & 1.522 & 26.928\% & 27.479\% & 45.593\% \\
Poisson & 1.173 & 1.355 & 31.343\% & 28.727\% & 39.930\% \\
Random Forest & 1.345 & 1.265 & 37.728\% & 28.311\% & 33.961\% \\
Gradient Boosting & 1.259 & 1.143 & 37.978\% & 29.672\% & 32.350\% \\
\bottomrule
\end{tabular}
}

\vspace{0.25cm}
\begin{block}{Lectura}
Con Ridge, Iran obtiene la mayor probabilidad de victoria y el marcador exacto más probable fue 1--1.
\end{block}
\end{frame}

%------------------------------------------------
\begin{frame}{Ajustes del modelo probabilístico}
\begin{block}{Dixon-Coles}
Se incorpora una corrección para mejorar la estimación de marcadores cortos, frecuentes en fútbol, como 0--0, 1--0, 0--1 y 1--1.

\[
P_{\text{ajustada}}(a,b)=P(a,b)\times \tau(a,b,\lambda_{\text{home}},\lambda_{\text{away}},\rho)
\]
\end{block}

\vspace{0.25cm}
\begin{block}{Contexto competitivo}
La expectativa de goles se ajusta según la situación del partido:
\[
\lambda_{\text{home-final}}=\lambda_{\text{home-base}}\times \text{factor}_{\text{home}}
\]

\[
\lambda_{\text{away-final}}=\lambda_{\text{away-base}}\times \text{factor}_{\text{away}}
\]
\end{block}

\vspace{0.2cm}
Este ajuste representa escenarios donde un equipo necesita atacar más, conservar un resultado o jugar de forma equilibrada.
\end{frame}

%------------------------------------------------
\begin{frame}{Comparación de modelos}
\small
\centering
\begin{tabular}{lcccc}
\toprule
Modelo & MAE prom. & RMSE prom. & $R^2$ prom. & Exactitud 1X2 \\
\midrule
Ridge & \textbf{0.975} & 1.315 & 0.279 & 0.624 \\
Regresión Lineal & \textbf{0.975} & 1.314 & 0.279 & 0.624 \\
Poisson & 0.977 & \textbf{1.307} & \textbf{0.288} & 0.625 \\
Random Forest & 0.985 & 1.316 & 0.278 & 0.623 \\
Gradient Boosting & 0.994 & 1.343 & 0.250 & \textbf{0.627} \\
\bottomrule
\end{tabular}

\vspace{0.35cm}
\begin{block}{Resultado de selección}
El modelo final elegido fue \textbf{Ridge}, por presentar el menor error absoluto promedio y una estructura estable para explicar la predicción.
\end{block}
\end{frame}

%------------------------------------------------
\begin{frame}{Lectura de los resultados}
\begin{itemize}
    \item Los errores promedio se ubican alrededor de un gol, una escala razonable para predicción futbolística.
    \item Ridge y Regresión Lineal tuvieron el menor MAE promedio.
    \item Poisson obtuvo el mejor RMSE promedio y el mayor $R^2$ promedio.
    \item Gradient Boosting alcanzó la mayor exactitud en el resultado 1X2.
    \item La selección final priorizó estabilidad, bajo error e interpretación.
\end{itemize}
\end{frame}

%------------------------------------------------
\begin{frame}{Ejemplo de análisis por partido}
\justifying
Para un partido seleccionado, el sistema compara todos los modelos y muestra:

\begin{itemize}
    \item Goles esperados de ambos equipos.
    \item Probabilidad de victoria local, empate y victoria visitante.
    \item Modelo ganador global.
    \item Marcadores exactos más probables.
    \item Efecto del contexto competitivo.
\end{itemize}

\vspace{0.25cm}
\begin{block}{Ejemplo evaluado}
Egypt vs Iran: el modelo Ridge estimó mayor probabilidad de victoria visitante, con el marcador 1--1 como marcador exacto más probable.
\end{block}
\end{frame}

%------------------------------------------------
\begin{frame}{Visualización comparativa}
\centering
\includegraphics[width=0.83\textwidth,height=0.72\textheight,keepaspectratio]{comparacion_modelos_match_65.png}
\end{frame}

%------------------------------------------------
\begin{frame}{Marcadores más probables}
\centering
\includegraphics[width=0.75\textwidth,height=0.72\textheight,keepaspectratio]{top10_match_65_ridge.png}
\end{frame}

%------------------------------------------------
\begin{frame}{Simulación de fase de grupos}
\justifying
Después de estimar los resultados de los partidos, se construyeron tablas de posiciones por grupo.

\vspace{0.25cm}
\begin{itemize}
    \item Se calcularon puntos, victorias, empates y derrotas.
    \item Se consideraron goles a favor, goles en contra y diferencia de goles.
    \item Se incorporó criterio de fair play para resolver escenarios de desempate.
    \item Se identificaron primeros, segundos y mejores terceros.
\end{itemize}

\vspace{0.25cm}
\begin{block}{Resultado global}
La simulación entrega \textbf{32 clasificados a 16avos}: 12 primeros, 12 segundos y 8 mejores terceros.
\end{block}
\end{frame}

%------------------------------------------------
\begin{frame}{Tablas predichas por grupo}
\centering
\includegraphics[width=0.88\textwidth,height=0.72\textheight,keepaspectratio]{panel_tablas_grupos.png}
\end{frame}

%------------------------------------------------
\begin{frame}{Clasificados y cruces simulados}
\centering
\includegraphics[width=0.48\textwidth,height=0.68\textheight,keepaspectratio]{panel_clasificados_16avos.png}
\hfill
\includegraphics[width=0.48\textwidth,height=0.68\textheight,keepaspectratio]{panel_cruces_16avos.png}
\end{frame}

%------------------------------------------------
\begin{frame}{Principales clasificados proyectados}
\small
\begin{columns}[T]
\column{0.5\textwidth}
\begin{block}{Primeros de grupo}
Mexico, Switzerland, Brazil, United States, Germany, Netherlands, Egypt, Spain, France, Argentina, Colombia y England.
\end{block}

\column{0.5\textwidth}
\begin{block}{Algunos cruces simulados}
Brazil vs Japan, Germany vs Sweden, France vs Turkey, Mexico vs Ecuador, Spain vs Austria y Argentina vs Cape Verde.
\end{block}
\end{columns}

\vspace{0.25cm}
Estos resultados representan una simulación basada en datos históricos, rankings y contexto competitivo; no son certezas deportivas.
\end{frame}

%------------------------------------------------
\begin{frame}{Conclusiones}
\begin{itemize}
    \item El enfoque permite pasar de datos históricos a predicciones interpretables de partidos.
    \item La combinación de regresión y Poisson permite estimar goles y convertirlos en probabilidades.
    \item Ridge fue seleccionado como modelo final por su bajo error promedio y estabilidad.
    \item El contexto competitivo mejora la lectura del partido porque incorpora la situación del torneo.
    \item La simulación final permite proyectar tablas, clasificados y cruces de eliminación.
\end{itemize}
\end{frame}

%------------------------------------------------
\begin{frame}{Mensaje final}
\centering
\Large
El proyecto no busca afirmar un resultado exacto, sino construir una herramienta analítica para estimar escenarios probables del Mundial 2026.

\vspace{0.5cm}
\normalsize
Gracias.
\end{frame}

\end{document}
